\clearpage
\pdfbookmark[1]{Coytier. 6 March 1813.}{anzeige_coytier_125_127}
\SourceLeaf{125}
\begin{center}
\par\noindent\hfill\rule{0.30\linewidth}{0.4pt}\hfill\mbox{}\par
Göttingische gelehrte Anzeigen. 6 March 1813.
\par\noindent\hfill\rule{0.30\linewidth}{0.4pt}\hfill\mbox{}\par
\end{center}

\textit{Mathematical investigations, or various questions that remain unsolved or whose solution} \textit{leaves something to be desired.} By P. L. \textsc{Coytier.} \textit{First memoir, containing} \textit{general observations on algebraic equations.} Paris 1813. \textit{Published by Eberhardt.} (16 pages in quarto.)\par

If we understand the wordy, bombastic preface correctly, the purpose of the present essay, and of others intended to follow it, is to use the weapons of synthesis to combat and overturn various results that, in the author's opinion, analysis has usurped. The work as a whole shows, however, that the author attaches notions of his own to the words analysis and synthesis: by the former he means mere blind juggling with symbols, by the latter a real, intuitive understanding. Such an undertaking certainly deserves praise if there is such blind juggling with symbols that presumes to call itself analysis. In most attempts of this kind, however, made on this side of the Rhine and beyond it, we find that the blindness lay in the attacker rather than in what was attacked. The contents, which we shall briefly describe, may serve to judge where the present work belongs.\par

The present work is chiefly directed against the factorization of algebraic equations into simple factors (this very expression already contains an error of which, admittedly, some other writers have also
\SourceLeaf{126}
been guilty). The author also often calls it the theorem of 1746, no doubt because \textsc{d’Alembert} first attempted a rigorous proof around that time. We expressed our judgement elsewhere 14 years ago that this proof of \textsc{d’Alembert}, like those of \textsc{Euler}, \textsc{Foncenex} and \textsc{Lagrange}, is by no means satisfactory; and, according to our conviction, we must of course judge the later work of \textsc{Laplace} and \textsc{Lagrange} on the same subject in the same way. Our author, however, does not enter into these proofs themselves at all: his remarks are of an entirely different kind and are directed against the theorem itself rather than against the proofs. He contends, for example (to keep only to the simplest case of a quadratic equation $xx-2ax+b=0$), that $x-[a+\surd(aa-b)]$ is not a simple factor at all, because $x-[a+\surd(aa-b)]=0$ is not an equation of the first degree at all, but a true quadratic equation. Whether the sign $+$ or the double sign $\pm$ is written before the radical quantity, he says, it always expresses both roots at once. Every equation, he maintains, actually represents the relation between two variable linear quantities and one constant linear quantity; the classification of equations and that of curves by degree must correspond most closely; the equation $x-[a+\surd(aa-b)]=0$ must be regarded as identical with the equation $xx-2ax+b=0$, and the former must therefore be counted as of the second degree just as much as the latter. These assertions make up the contents of the work and show us only the confusion of the author's notions. Ordinary algebra knows no variable quantities at all, only unknown and known ones. The radical sign $\surd$ actually has a double meaning in mathematical notation (and this is admittedly a small imperfection): $\surd A$ is to be defined either as \textit{a} quantity whose square $=A$, or as \textit{the} positive quantity whose square $=A$, insofar as $A$ is positive. It is up to the analyst how he wishes to use the sign; a thoughtful, careful analyst will always be clearly aware of this and will always express himself so that the reader too is left in no doubt whether the sign is used in an indeterminate or a determinate sense. As we already indicated above, \textit{equations} are not factorized at all, but functions of \textit{one} variable quantity. Thus it is not the equation $xx-2ax+b=0$ but the function $xx-2ax+b$ that is resolved into the factors $x-[a+\surd(aa-b)]$, $x-[a-\surd(aa-b)]$, and inso\SourceLeaf{127}%
far as only $x$ is regarded as variable in them, they are \textit{called} factors of the first degree. The author is welcome, when he regards the equation $xx-2ax+b=0$ as the expression of a curve and $a$ or $b$ as variable, to call the factor $x-[a+\surd(aa-b)]$ an expression of the second degree. In doing so he disputes no \textit{truth}, but only the appropriateness of a \textit{name} in a case where it is not used and never has been used.\par

Most mathematical readers of these pages will, of course, find these discussions unnecessary. We nevertheless believe that we must develop the author's notions a little, so that no one is led by the work's title to expect new insights for which the author seems to lack the very rudiments.\par
\par\noindent\hfill\rule{0.30\linewidth}{0.4pt}\hfill\mbox{}\par
